\section{Artifact Appendix}

\subsection{Abstract}
Our artifact provides the C++ implementation of fuzzy private set
intersection (FPSI) under the unique-cell and unique-block assumptions,
including normal and prefix-optimized protocols for $L_\infty$, $L_1$,
and $L_2$ distances. It includes build scripts, generated-input correctness
tests, performance benchmarks, and paper-reference data. Evaluation produces
raw logs, Markdown summaries, and SVG runtime comparisons with the
\emph{Ours}/\emph{Ours-Px} entries in Tables~2 and~3. Our code is MIT-licensed;
dependencies retain their respective licenses.

\subsection{Description \& Requirements}
The root \path{README.md} is the entry point. Implementation modules are
under \path{code/fpsi/}; \path{scripts/} contains categorized build,
reproduction, and data tools. The \path{claims/} directory maps each claim to its commands,
expected outputs, and paper-reference CSV files. Generated results are
written to \path{artifact-results/}.

\subsubsection{Security, privacy, and ethical concerns}
Benchmarks use synthetic, non-sensitive points and execute locally. Installation
downloads dependencies and may install system packages with root or
\texttt{sudo}; benchmarks do not require either. Large cases can exhaust RAM;
use an otherwise idle host with the memory budget below. This is a research
prototype, not a production security implementation: the shared-output OPRF
uses fixed benchmark keys and preconstructed key-OT inputs. Tests validate
functional outputs and performance, not deployment security.

\subsubsection{How to access}
Download: \ArtifactURL. Permanent archive: \ArtifactDOI.
Evaluated version: \ArtifactVersion. Contact: \ArtifactContact.

\subsubsection{Hardware dependencies}
Use an x86-64 CPU supporting AES, PCLMUL, SSE2, and SSE4.1;
ARM64 and Apple Silicon are unsupported. No GPU or second host is needed:
both protocol parties run as local threads in one process. Build and basic
validation require 16\,GiB RAM. For full evaluation, recommend 256\,GiB RAM;
observed peak process RSS is approximately 63\,GiB for unique-cell and
199\,GiB for unique-block. Budgets below use an AMD EPYC 9554 host and
one trial; slower hosts require more time. Leave at least 10\,GiB free disk
space for dependencies, build products, and results.

\subsubsection{Software dependencies}
The reference platform is Ubuntu 24.04; GCC 11 or newer with C++20 is
required. Build scripts fetch pinned dependencies.
Bash, CMake, Git, Python~3, and the packages listed in the root README are
required. Internet access is needed for a fresh build. Result processing and
SVG plotting use Python's standard library and our modules only; Matplotlib
is not required. Docker is the recommended deployment method.

\subsubsection{Benchmarks}
No external datasets or models are required. Benchmark inputs are generated
unsigned-integer point sets satisfying the selected one-sided assumption.
Performance scripts explicitly plant 16 matches (\texttt{-inter 16}). Correctness tests generate
edge-case datasets and reference answers.

\subsection{Set Up}
\subsubsection{Installation}
Obtain the source archive, extract it, and run from its root:
\begin{lstlisting}
./scripts/build/preflight.sh
./scripts/build/run.sh
\end{lstlisting}
Allow about 5 person-minutes and 10--20 compute-minutes for a fresh build on
an 8-core machine. Preflight checks the host without modifying it; on a fresh
host, install any missing prerequisites reported and rerun it. A successful
build creates \path{code/build/fpsi} and reports \texttt{Build complete}.
If system packages are already installed, use
\texttt{FPSI\_SKIP\_SYSTEM\_PACKAGES=1} with the build command to avoid package
installation. Dependency versions and troubleshooting are in the root README.

\subsubsection{Basic test}
Run \path{./scripts/reproduction/run.sh} \texttt{--quick}
(about 1 person-minute, under 2 compute-minutes, 16\,GiB RAM).
The smoke test must report
\texttt{PASS: 6 protocol cases and 1 parameter guard}.
Inspect \path{artifact-results/quick.txt} for the raw output and
\path{summary.md} for measurements. This is an installation check,
not the complete correctness experiment.

\subsection{Evaluation Workflow}
\subsubsection{Major claims}
\begin{compactitem}
\item[(C1)] The implemented protocols correctly compute FPSI by returning all
sender elements that are close to the receiver set (E1).
\item[(C2)] The receiver-sided unique-cell $L_\infty$ implementations
reproduce the runtime and communication trends of the
\emph{Ours}/\emph{Ours-Px} entries in Table~2 (E2).
\item[(C3)] The receiver-sided unique-block implementations reproduce
the corresponding trends for $L_\infty$, $L_1$, and $L_2$ in Table~3 (E3).
\end{compactitem}
These experiments do not rerun competing baselines or establish security
proofs. The performance claims concern our table entries, not independent
reproduction of every baseline comparison.

Reported seconds use each protocol's original timed interval, not process
wall time: synthetic input generation, sockets, and mode-specific local
preprocessing are excluded; OPRF, OKVS, MPC, final transfer, and enabled
correctness checks are timed. The exact per-mode split is documented in
\path{claims/paper-comparison.md}. Communication is the sum sent and received
by one local endpoint, divided by $2^{20}$ bytes (labeled MB in the tables).
Default runs do not apply network shaping.

\subsubsection{Experiments}
Run E1--E3 sequentially from the artifact root after installation. Budgets
exclude build time and use the default \texttt{TRIALS=1}. Each wrapper exits
nonzero on validation failure and preserves diagnostic logs. Increase
\texttt{TRIALS}, e.g., \texttt{TRIALS=3} \texttt{bash} \path{claims/claim2/run.sh}, to obtain
internal-trial mean measurements; budget proportionally more compute time.

\paragraph{E1: Correctness (C1).}
\emph{5 person-minutes + 3--5 compute-minutes; 16\,GiB RAM.}
\textbf{Preparation:} No input download is needed.
\textbf{Execution:} Run \texttt{bash} \path{claims/claim1/run.sh}.
\textbf{Results:} The wrapper runs smoke checks and 240 applicable
dataset/configuration executions covering every implemented protocol variant.
Cases include zero/full
intersections, many-to-one matches, distances $\delta-1$, $\delta$, and
$\delta+1$, and cell/block-boundary matches. An independent Python
integer-distance oracle checks every recovered sender point, not just the
match count. The final report must include
\texttt{240 executions across 16 modes; 0 failures}.
All 240 rows in \path{artifact-results/claim1/}\allowbreak
\path{boundary/runs/summary.csv} must be \texttt{PASS}; per-case inputs,
expected answers, outputs, and logs remain available.

\paragraph{E2: Unique-cell (C2, Table 2).}
\emph{10 person-minutes + 3--3.5 compute-hours; 80\,GiB RAM recommended.}
\textbf{Preparation:} Reserve an otherwise idle host.
\textbf{Execution:} Run \texttt{bash} \path{claims/claim2/run.sh}.
Both modes cover
$n=m\in\{2^8,2^{12},2^{16}\}$, $d\in\{2,4,6\}$, and
$\delta\in\{32,64,128,256,512\}$: 90 configurations.
\textbf{Results:} Under \path{artifact-results/claim2/}, expect 90 rows
in both \path{summary.md} and \path{paper-comparison.md}, and nine
runtime SVGs in \path{paper-plots/}, one per $(n,d)$ pair. The terminal reports
\texttt{Claim 2 full complete}. For smaller runs, use the shared partial
reproduction or mini benchmark described in Notes on Reusability.

\paragraph{E3: Unique-block (C3, Table 3).}
\emph{10 person-minutes; 2--2.5 hours compute; 256\,GiB RAM recommended.}
\textbf{Preparation:} Reserve the full memory budget, especially for normal
mode at $d=6$.
\textbf{Execution:} Run \texttt{bash} \path{claims/claim3/run.sh}.
This evaluates $n=m=2^{12}$, the same dimensions and thresholds as E2,
all three metrics, and both modes: 90 configurations.
\textbf{Results:} Under \path{artifact-results/claim3/}, expect 90 rows
in each summary/comparison Markdown table and nine runtime SVGs, one per metric/dimension pair.
The terminal reports \texttt{Claim 3 complete}.

For E2/E3, each trial must report \texttt{Total 16/16 matches found!}.
Wrappers reject missing or mismatched configurations. Open each
\path{paper-comparison.md} to inspect measured and paper runtime curves
against equally spaced $\delta$ categories, in seconds with automatically
scaled linear time axes.
Communication should be stable for fixed parameters; absolute runtime varies
with hardware. Assess scaling across thresholds, dimensions, and set sizes;
do not require exact times or strictly increasing adjacent samples.
The plots support manual interpretation, not an automatic trend-pass verdict.

\subsection{Notes on Reusability}
For lower-cost exploration, use
\path{./scripts/reproduction/run.sh} \texttt{--partial}: 80 original-paper
configurations at $n=2^{12},d\in\{2,4\}$, about 25--35 minutes and 35\,GiB
peak RSS (64\,GiB RAM recommended). The \texttt{--mini} profile uses the same
matrix at $n=2^{10}$: about 6--10 minutes and 21\,GiB peak RSS (32\,GiB RAM).
Mini is not a reproduction of paper set sizes. These profiles produce logs
and summaries; use E2/E3 for automatic paper-comparison plots.

Custom data can be supplied using \texttt{-i <directory>} with
\path{sender_data.txt} and \path{recver_data.txt}: one point per line,
whitespace-separated unsigned coordinates, equal set sizes and dimensions,
no duplicates within either set, and the selected input assumption satisfied.
Matching sender points are written to \path{output.txt} in that directory.
Without \texttt{-i}, simulated inputs are generated. The reusable boundary
generator and modular protocol/primitive code are documented in
\path{claims/claim1/claim.md} and \path{code/README.md}.

\subsection{Version}
Based on the LaTeX template for Artifact Evaluation V20220926.
